\section{Mathematical Evaluation Methodology}
\label{sec:methodology}

We evaluate our routing engine using three mathematical approaches: multivariate Ordinary Least Squares (OLS) regression, multi-factor Analysis of Variance (ANOVA), and K-Means clustering. 
We execute 1500 random route queries across the Moscow metropolitan road network to collect performance metrics. 

We apply multivariate OLS regression to isolate execution time components. 
Our regression model estimates query latency $T_{\text{ms}}$ based on priority queue operation costs, relaxation steps, and cache miss rates. 
The OLS model measures coefficients for each queue layout to quantify microarchitectural overheads. 

We perform multi-factor ANOVA testing to establish the significance of priority queue selection on performance variance. 
We track two factors: the search algorithm (Dijkstra, A*, or ALT) and the priority queue type. 
ANOVA checks the hypothesis that queue selection exerts a uniform influence on execution latency across different routing paradigms. 
We measure latency jitter using variance, skewness, and kurtosis to evaluate stability under OS scheduler noise.

We classify the 1500 benchmark journeys using K-Means clustering to analyze scaling patterns. 
The clustering algorithm leverages journey distance, edge count, and node exploration counts as features. 
We identify three distinct routing profiles:
\begin{itemize}
\item \textbf{Localized Profile}: Cluster 2, mean distance 12.3\,km. Graph traversal overhead is small. Initialization and fast-clear costs dominate the execution profile.
\item \textbf{Regional Profile}: Cluster 0, mean distance 28.4\,km. Traversal cost and queue operations exhibit equal latency footprints. 
\item \textbf{Megalopolis Profile}: Cluster 1, mean distance 43.7\,km. The priority queue hot loop determines execution time.
\end{itemize}
This classification allows us to observe how priority queue performance scales from local queries to long-distance paths.
